% =============================================================================
%  cours_sql.tex - SQL : des bases a la preparation d'entretien technique
%  Syntaxe de reference : PostgreSQL (differences MySQL notees quand utiles).
%
%  Compilation (deux passes pour la table des matieres) :
%     pdflatex cours_sql.tex
%     pdflatex cours_sql.tex
% =============================================================================
\documentclass[11pt,a4paper]{article}

\usepackage[utf8]{inputenc}
\usepackage[T1]{fontenc}
\usepackage[french]{babel}
\usepackage[margin=2.2cm]{geometry}
\usepackage[table]{xcolor}
\usepackage{listings}
\usepackage{tcolorbox}
\usepackage{enumitem}
\usepackage{titlesec}
\usepackage{longtable}
\usepackage{array}
\usepackage{fancyhdr}
\usepackage{hyperref}

% --- Couleurs ---------------------------------------------------------------
\definecolor{primary}{HTML}{1B4B5A}    % bleu petrole (theme SQL)
\definecolor{accent}{HTML}{2E86DE}     % bleu accent
\definecolor{codebg}{HTML}{F4F7FB}
\definecolor{codekw}{HTML}{1B4B5A}
\definecolor{codestr}{HTML}{B7472A}
\definecolor{codecom}{HTML}{7F8C99}
\definecolor{rulecol}{HTML}{D5DEE8}
\definecolor{warnbg}{HTML}{FFF4E6}
\definecolor{warnframe}{HTML}{E67E22}

\hypersetup{colorlinks=true, linkcolor=accent, urlcolor=accent, pdftitle={Cours SQL - des bases a l'entretien technique}}

% --- En-tete / pied de page -------------------------------------------------
\pagestyle{fancy}
\fancyhf{}
\fancyhead[L]{\small\color{primary}Cours SQL -- des bases a l'entretien technique}
\fancyhead[R]{\small\color{codecom}\thepage}
\renewcommand{\headrulewidth}{0.4pt}

% --- Titres -------------------------------------------------------------
\titleformat{\section}{\Large\bfseries\color{primary}}{\thesection.}{0.6em}{}
\titleformat{\subsection}{\large\bfseries\color{accent}}{\thesubsection}{0.6em}{}

% --- Style du code SQL -------------------------------------------------------
\lstdefinestyle{sql}{
  language=SQL,
  backgroundcolor=\color{codebg},
  basicstyle=\ttfamily\small,
  keywordstyle=\color{codekw}\bfseries,
  stringstyle=\color{codestr},
  commentstyle=\color{codecom}\itshape,
  numbers=none,
  frame=single,
  rulecolor=\color{rulecol},
  breaklines=true,
  showstringspaces=false,
  columns=fullflexible,
  keepspaces=true,
  aboveskip=8pt, belowskip=8pt,
  xleftmargin=6pt, xrightmargin=6pt,
  morekeywords={ILIKE, RETURNING, WINDOW, PARTITION, OVER, LATERAL, WITH,
    RECURSIVE, UPSERT, ONCONFLICT, CONFLICT, EXCLUDED, FILTER, CASCADE,
    RESTRICT, SERIAL, BIGSERIAL, TEXT, JSONB, TRUE, FALSE},
  literate=%
    {é}{{\'e}}1 {è}{{\`e}}1 {ê}{{\^e}}1 {à}{{\`a}}1 {ô}{{\^o}}1
    {î}{{\^i}}1 {ç}{{\c c}}1 {É}{{\'E}}1 {È}{{\`E}}1 {ù}{{\`u}}1 {û}{{\^u}}1,
}
\lstset{style=sql}

% --- Boites -----------------------------------------------------------------
\newtcolorbox{infobox}[1]{colback=codebg,colframe=accent,title=#1,
  fonttitle=\bfseries\color{white},boxrule=0.5pt,arc=2pt,left=6pt,right=6pt,top=4pt,bottom=4pt}

\newtcolorbox{piege}{colback=warnbg,colframe=warnframe,title=\textbf{Piege courant},
  fonttitle=\bfseries\color{white},boxrule=0.5pt,arc=2pt,left=6pt,right=6pt,top=4pt,bottom=4pt}

% =============================================================================
\begin{document}

% --- Page de titre ----------------------------------------------------------
\begin{titlepage}
\centering
\vspace*{3cm}
{\color{primary}\Huge\bfseries SQL\par}
\vspace{0.6cm}
{\color{accent}\LARGE Des bases a la preparation d'entretien technique\par}
\vspace{1.2cm}
{\large Requetes, jointures, agregations, fonctions fenetrees,\\
CTE et index -- avec la syntaxe PostgreSQL.\par}
\vfill
{\color{codecom}\large Cours a imprimer -- support d'entrainement\par}
\vspace{2cm}
\end{titlepage}

\tableofcontents
\newpage

% =============================================================================
\section{Introduction}
% =============================================================================
SQL (\emph{Structured Query Language}) est le langage utilise pour interroger et
manipuler des donnees stockees dans une base de donnees \emph{relationnelle} :
des tables composees de lignes et de colonnes, reliees entre elles par des
\emph{cles}. Contrairement a un langage comme Python, SQL est \emph{declaratif} :
on decrit \emph{quel} resultat on veut, pas \emph{comment} l'obtenir -- c'est le
\emph{moteur} de la base (le \emph{query planner}) qui choisit la meilleure
strategie d'execution.

Les exemples de ce cours utilisent la syntaxe \textbf{PostgreSQL}. L'essentiel
est du SQL standard (norme ANSI) et fonctionne tel quel sur la plupart des
moteurs (MySQL, SQL Server, SQLite) ; les differences notables sont signalees
dans des encadres.

\begin{infobox}{Les grandes familles de commandes}
\textbf{DDL} (\emph{Data Definition Language}) : \texttt{CREATE}, \texttt{ALTER},
\texttt{DROP} -- definissent la structure. \textbf{DML} (\emph{Data Manipulation
Language}) : \texttt{SELECT}, \texttt{INSERT}, \texttt{UPDATE}, \texttt{DELETE}
-- manipulent les donnees. \textbf{DCL} (\emph{Data Control Language}) :
\texttt{GRANT}, \texttt{REVOKE} -- gerent les droits d'acces.
\end{infobox}

% =============================================================================
\section{Types de donnees}
% =============================================================================
Chaque colonne d'une table a un \emph{type} fixe, choisi a la creation. Le type
determine ce que le SGBD accepte d'y stocker, l'espace occupe sur disque, et les
operations autorisees (on ne peut pas faire de moyenne sur du texte). Bien
choisir ses types des le depart evite des migrations couteuses plus tard.

\begin{longtable}{>{\raggedright\arraybackslash}p{3.6cm}
                  >{\raggedright\arraybackslash}p{9.2cm}}
\rowcolor{primary}
\textbf{\color{white}Categorie} & \textbf{\color{white}Types courants (PostgreSQL)} \\
\hline
Entiers & \texttt{SMALLINT} (16 bits), \texttt{INTEGER} (32 bits),
\texttt{BIGINT} (64 bits). \texttt{SERIAL}/\texttt{BIGSERIAL} : entier
auto-incremente, pratique pour une cle primaire. \\
Nombres decimaux & \texttt{NUMERIC(precision, echelle)} : precision exacte,
pour de l'argent. \texttt{REAL}, \texttt{DOUBLE PRECISION} : virgule flottante
approximative, pour du calcul scientifique. \\
Texte & \texttt{TEXT} : longueur illimitee. \texttt{VARCHAR(n)} : longueur
limitee a \texttt{n}. \texttt{CHAR(n)} : longueur fixe, complete avec des
espaces (rarement utile). \\
Booleen & \texttt{BOOLEAN} : \texttt{TRUE}, \texttt{FALSE}, ou \texttt{NULL}
(trois etats, pas deux). \\
Date et heure & \texttt{DATE} (jour seul), \texttt{TIME} (heure seule),
\texttt{TIMESTAMP} (date + heure, sans fuseau), \texttt{TIMESTAMPTZ}
(date + heure, avec fuseau -- recommande des que l'app a des utilisateurs
dans plusieurs fuseaux). \\
Identifiants & \texttt{UUID} : identifiant unique global, utile pour eviter
des ID predictibles ou fusionner des donnees de plusieurs sources. \\
Structures & \texttt{JSON}/\texttt{JSONB} : donnees semi-structurees (JSONB
est stocke en binaire, indexable, generalement prefere). \texttt{ARRAY} :
tableau d'un type donne (\texttt{INTEGER[]}, \texttt{TEXT[]}). \\
\end{longtable}

\begin{lstlisting}
CREATE TABLE produits (
    id           BIGSERIAL PRIMARY KEY,
    nom          VARCHAR(120)   NOT NULL,
    description  TEXT,
    prix         NUMERIC(10, 2) NOT NULL,      -- 10 chiffres, dont 2 apres la virgule
    en_stock     BOOLEAN        DEFAULT TRUE,
    tags         TEXT[],                        -- tableau de textes
    attributs    JSONB,                          -- donnees libres, indexables
    cree_le      TIMESTAMPTZ    DEFAULT NOW()
);

-- Conversion explicite entre types avec CAST ou ::
SELECT CAST('42' AS INTEGER);
SELECT '42'::INTEGER;              -- syntaxe raccourcie propre a PostgreSQL
\end{lstlisting}

\begin{infobox}{NUMERIC vs FLOAT}
\texttt{NUMERIC} stocke une valeur \emph{exacte} (aucune erreur d'arrondi) mais
est plus lent ; c'est le choix par defaut pour de l'argent. \texttt{REAL}/
\texttt{DOUBLE PRECISION} sont des flottants binaires, plus rapides mais avec
des approximations (\texttt{0.1 + 0.2} peut ne pas donner exactement
\texttt{0.3}) -- a reserver aux calculs ou une petite imprecision est
acceptable.
\end{infobox}

\begin{piege}
\texttt{VARCHAR(n)} sans \texttt{n} et \texttt{TEXT} sont strictement
equivalents en PostgreSQL (aucune difference de performance) : la limite de
longueur n'apporte qu'une validation, pas un gain de vitesse. Ne pas confondre
avec MySQL, ou \texttt{VARCHAR} et \texttt{TEXT} sont stockes differemment.
\end{piege}

% =============================================================================
\section{Creer et peupler des tables}
% =============================================================================
Une table est definie par ses colonnes, chacune avec un \emph{type} et
d'eventuelles \emph{contraintes}. Les contraintes empechent les donnees
invalides d'entrer dans la base -- c'est la premiere ligne de defense de
l'integrite des donnees, avant meme le code applicatif.

\begin{lstlisting}
CREATE TABLE clients (
    id          SERIAL PRIMARY KEY,          -- entier auto-incremente
    nom         TEXT NOT NULL,
    email       TEXT UNIQUE NOT NULL,
    age         INTEGER CHECK (age >= 0),
    cree_le     TIMESTAMP DEFAULT NOW()
);

CREATE TABLE commandes (
    id          SERIAL PRIMARY KEY,
    client_id   INTEGER NOT NULL REFERENCES clients(id),  -- cle etrangere
    montant     NUMERIC(10, 2) NOT NULL,
    statut      TEXT DEFAULT 'en_attente'
);

INSERT INTO clients (nom, email, age)
VALUES ('Alice', 'alice@mail.com', 30),
       ('Bob',   'bob@mail.com',   25);

UPDATE clients SET age = 31 WHERE nom = 'Alice';

DELETE FROM clients WHERE email = 'bob@mail.com';
\end{lstlisting}

\begin{piege}
\texttt{DELETE FROM table} sans clause \texttt{WHERE} supprime \textbf{toutes}
les lignes. \texttt{UPDATE table SET ...} sans \texttt{WHERE} modifie
\textbf{toutes} les lignes. Toujours verifier la clause \texttt{WHERE} avant
d'executer, surtout en production -- idealement, tester d'abord avec un
\texttt{SELECT} utilisant la meme condition.
\end{piege}

% =============================================================================
\section{Selectionner des donnees}
% =============================================================================
\texttt{SELECT} est la commande la plus utilisee : elle extrait des lignes
d'une ou plusieurs tables. L'ordre \emph{logique} d'execution d'une requete
n'est pas l'ordre dans lequel on l'ecrit -- comprendre cet ordre evite
beaucoup d'erreurs (par exemple, pourquoi un alias defini dans \texttt{SELECT}
n'est pas utilisable dans \texttt{WHERE}).

\begin{infobox}{Ordre logique d'execution}
\texttt{FROM} $\rightarrow$ \texttt{JOIN} $\rightarrow$ \texttt{WHERE}
$\rightarrow$ \texttt{GROUP BY} $\rightarrow$ \texttt{HAVING} $\rightarrow$
\texttt{SELECT} $\rightarrow$ \texttt{DISTINCT} $\rightarrow$ \texttt{ORDER BY}
$\rightarrow$ \texttt{LIMIT}. \texttt{WHERE} filtre les lignes \emph{avant}
regroupement ; \texttt{HAVING} filtre les groupes \emph{apres}.
\end{infobox}

\begin{lstlisting}
SELECT nom, age
FROM clients
WHERE age >= 18 AND nom ILIKE 'a%'      -- ILIKE = LIKE insensible a la casse
ORDER BY age DESC
LIMIT 10;

-- Colonnes calculees et alias
SELECT nom, age, age >= 18 AS majeur
FROM clients;

-- Valeurs distinctes
SELECT DISTINCT statut FROM commandes;

-- Tester une valeur NULL : jamais avec = ou !=
SELECT * FROM clients WHERE email IS NOT NULL;

-- Plage de valeurs et liste de valeurs
SELECT * FROM clients WHERE age BETWEEN 18 AND 65;
SELECT * FROM commandes WHERE statut IN ('en_attente', 'expediee');
\end{lstlisting}

\begin{piege}
\texttt{NULL} represente une valeur \emph{inconnue}, pas zero ni chaine vide.
\texttt{colonne = NULL} ne renvoie jamais \texttt{TRUE} (meme si la colonne est
vraiment \texttt{NULL}) : il faut \texttt{IS NULL} ou \texttt{IS NOT NULL}.
\end{piege}

\subsection{Astuces de manipulation}
Au-dela du simple filtrage, quelques fonctions et operateurs reviennent sans
cesse pour \emph{extraire}, \emph{nettoyer} ou \emph{transformer} une colonne
directement dans la requete -- sans avoir a rapatrier les donnees brutes pour
les retravailler cote application.

\begin{lstlisting}
-- Debut / fin d'une chaine : LEFT et RIGHT prennent un nombre de caracteres
SELECT LEFT(nom, 3)   AS debut,     -- 3 premiers caracteres
       RIGHT(nom, 3)  AS fin        -- 3 derniers caracteres
FROM clients;

-- SUBSTRING : extraire depuis une position, sur une longueur donnee
SELECT SUBSTRING(email FROM 1 FOR 5)              AS debut_email,
       SUBSTRING(email FROM POSITION('@' IN email) + 1) AS domaine
FROM clients;

-- LENGTH : longueur d'une chaine (utile combine a LEFT/RIGHT pour du dynamique)
SELECT nom, LENGTH(nom) AS longueur FROM clients;

-- Nettoyage : equivalents SQL de strip/lstrip/rstrip en Python
SELECT TRIM(nom), LTRIM(nom), RTRIM(nom) FROM clients;
SELECT TRIM(BOTH '#' FROM '##titre##');            -- 'titre'

-- Casse et concatenation
SELECT UPPER(nom), LOWER(email) FROM clients;
SELECT nom || ' <' || email || '>' AS contact FROM clients;   -- || = concatenation
SELECT CONCAT(nom, ' <', email, '>') AS contact FROM clients; -- equivalent, gere les NULL

-- Remplacement
SELECT REPLACE(email, '@mail.com', '@nouveau.com') FROM clients;

-- COALESCE : premiere valeur non-NULL dans une liste -- LA methode pour
-- gerer les valeurs manquantes sans passer par un CASE WHEN
SELECT nom, COALESCE(telephone, 'non renseigne') AS tel FROM clients;

-- NULLIF : renvoie NULL si les deux valeurs sont egales (utile pour eviter
-- une division par zero sans passer par un CASE WHEN)
SELECT total / NULLIF(quantite, 0) AS prix_unitaire FROM commandes;

-- CASE WHEN : la seule structure conditionnelle de SQL, comme un if/elif/else
SELECT nom, age,
       CASE
           WHEN age < 18 THEN 'mineur'
           WHEN age < 65 THEN 'adulte'
           ELSE 'senior'
       END AS categorie
FROM clients;
\end{lstlisting}

\begin{infobox}{Alignement des positions}
En SQL, les positions de chaine commencent a \texttt{1} (pas \texttt{0} comme
en Python). \texttt{LEFT(s, n)} et \texttt{RIGHT(s, n)} avec un \texttt{n}
negatif retirent \texttt{n} caracteres depuis l'autre bout
(\texttt{LEFT('bonjour', -2)} donne \texttt{'bonj'}) -- pratique mais facile a
oublier.
\end{infobox}

\begin{piege}
\texttt{COALESCE} et \texttt{CASE WHEN} s'arretent a la premiere condition
vraie / premiere valeur non-\texttt{NULL} : l'ordre des arguments compte.
Utiliser \texttt{||} pour concatener avec une valeur \texttt{NULL} propage le
\texttt{NULL} a toute la chaine resultante (\texttt{'a' || NULL} donne
\texttt{NULL}, pas \texttt{'a'}) ; \texttt{CONCAT} n'a pas ce probleme, il
traite \texttt{NULL} comme une chaine vide.
\end{piege}

% =============================================================================
\section{Jointures}
% =============================================================================
Une jointure combine les lignes de deux tables selon une condition de
correspondance (generalement une cle etrangere). C'est le mecanisme qui
justifie de \emph{normaliser} les donnees (les repartir dans plusieurs tables
sans duplication) plutot que de tout stocker dans une seule table plate.

\begin{lstlisting}
-- INNER JOIN : seulement les lignes qui correspondent des deux cotes
SELECT c.nom, o.montant
FROM clients c
INNER JOIN commandes o ON o.client_id = c.id;

-- LEFT JOIN : toutes les lignes de gauche, meme sans correspondance a droite
-- (colonnes de droite a NULL si pas de commande)
SELECT c.nom, o.montant
FROM clients c
LEFT JOIN commandes o ON o.client_id = c.id;

-- Clients SANS aucune commande (anti-jointure classique)
SELECT c.nom
FROM clients c
LEFT JOIN commandes o ON o.client_id = c.id
WHERE o.id IS NULL;

-- FULL OUTER JOIN : union des deux cotes, correspondance ou non
SELECT c.nom, o.montant
FROM clients c
FULL OUTER JOIN commandes o ON o.client_id = c.id;

-- Jointure d'une table sur elle-meme (self-join) : ex. trouver les paires
-- d'employes qui partagent le meme manager
SELECT e1.nom AS employe, e2.nom AS collegue
FROM employes e1
JOIN employes e2 ON e1.manager_id = e2.manager_id AND e1.id < e2.id;
\end{lstlisting}

\begin{infobox}{Quelle jointure choisir ?}
\texttt{INNER JOIN} : uniquement les correspondances (le plus courant).
\texttt{LEFT JOIN} : tout de la table de gauche, complete par \texttt{NULL} si
pas de correspondance -- utile pour « trouver ce qui manque ». \texttt{FULL
OUTER JOIN} : rare, sert a comparer deux ensembles en gardant tout des deux
cotes.
\end{infobox}

\begin{piege}
Filtrer sur la table de droite d'un \texttt{LEFT JOIN} dans le \texttt{WHERE}
(au lieu du \texttt{ON}) transforme silencieusement le \texttt{LEFT JOIN} en
\texttt{INNER JOIN} : \texttt{WHERE o.statut = 'payee'} elimine les clients
sans commande alors qu'on voulait les garder. La condition sur la table jointe
doit rester dans le \texttt{ON} si on veut conserver les lignes sans
correspondance.
\end{piege}

% =============================================================================
\section{Agregations et regroupement}
% =============================================================================
Les fonctions d'agregation (\texttt{COUNT}, \texttt{SUM}, \texttt{AVG},
\texttt{MIN}, \texttt{MAX}) resument plusieurs lignes en une seule valeur.
\texttt{GROUP BY} les applique \emph{par groupe} plutot que sur toute la table :
chaque valeur distincte des colonnes groupees devient une ligne de resultat.
Toute colonne du \texttt{SELECT} qui n'est pas agregee doit apparaitre dans le
\texttt{GROUP BY}.

\begin{lstlisting}
-- Nombre et montant total de commandes par client
SELECT client_id, COUNT(*) AS nb_commandes, SUM(montant) AS total
FROM commandes
GROUP BY client_id;

-- HAVING filtre les groupes (apres agregation) ; WHERE filtre les lignes (avant)
SELECT client_id, SUM(montant) AS total
FROM commandes
WHERE statut = 'payee'          -- filtre les lignes avant regroupement
GROUP BY client_id
HAVING SUM(montant) > 100;      -- filtre les groupes apres regroupement

-- COUNT(*) compte les lignes ; COUNT(colonne) ignore les NULL de cette colonne
SELECT COUNT(*) AS total_lignes, COUNT(email) AS emails_renseignes
FROM clients;
\end{lstlisting}

\begin{piege}
\texttt{WHERE} ne peut pas utiliser un alias defini dans \texttt{SELECT}, ni
filtrer sur le resultat d'une agregation (\texttt{WHERE COUNT(*) > 1} est
invalide) : l'agregation n'existe pas encore a ce stade de l'execution logique.
C'est justement le role de \texttt{HAVING}.
\end{piege}

% =============================================================================
\section{Sous-requetes}
% =============================================================================
Une sous-requete est une requete imbriquee dans une autre. Elle peut apparaitre
dans le \texttt{WHERE} (filtrer), dans le \texttt{FROM} (comme une table
temporaire), ou dans le \texttt{SELECT} (une valeur calculee par ligne).

\begin{lstlisting}
-- Sous-requete scalaire : clients dont l'age depasse la moyenne
SELECT nom, age
FROM clients
WHERE age > (SELECT AVG(age) FROM clients);

-- Sous-requete avec IN : clients ayant au moins une commande payee
SELECT nom FROM clients
WHERE id IN (SELECT client_id FROM commandes WHERE statut = 'payee');

-- EXISTS : souvent plus rapide que IN sur de gros volumes,
-- s'arrete des la premiere ligne trouvee
SELECT nom FROM clients c
WHERE EXISTS (
    SELECT 1 FROM commandes o
    WHERE o.client_id = c.id AND o.statut = 'payee'
);

-- Sous-requete correlee dans le SELECT : montant total par client, en ligne
SELECT nom,
       (SELECT SUM(montant) FROM commandes o WHERE o.client_id = c.id) AS total
FROM clients c;
\end{lstlisting}

% =============================================================================
\section{CTE -- Common Table Expressions}
% =============================================================================
Une CTE (clause \texttt{WITH}) nomme une sous-requete pour la reutiliser et
rendre une requete complexe lisible, comme on decoupe un long calcul en
variables intermediaires. Une CTE \emph{recursive} permet de parcourir des
structures hierarchiques (arbre, graphe) que \texttt{SQL} ne sait pas
traverser nativement autrement.

\begin{lstlisting}
-- CTE simple : lisibilite d'une requete en plusieurs etapes
WITH ventes_par_client AS (
    SELECT client_id, SUM(montant) AS total
    FROM commandes
    GROUP BY client_id
)
SELECT c.nom, v.total
FROM clients c
JOIN ventes_par_client v ON v.client_id = c.id
WHERE v.total > 500;

-- CTE recursive : remonter toute la hierarchie d'un employe (ex. organigramme)
WITH RECURSIVE hierarchie AS (
    -- cas de base : l'employe de depart
    SELECT id, nom, manager_id, 1 AS niveau
    FROM employes
    WHERE id = 42

    UNION ALL

    -- cas recursif : le manager du precedent, niveau par niveau
    SELECT e.id, e.nom, e.manager_id, h.niveau + 1
    FROM employes e
    JOIN hierarchie h ON e.id = h.manager_id
)
SELECT * FROM hierarchie;
\end{lstlisting}

\begin{infobox}{Quand utiliser une CTE}
Des qu'une sous-requete est reutilisee plusieurs fois, ou que la requete
devient difficile a lire d'un seul bloc. La CTE recursive est le seul moyen
standard de parcourir un nombre variable de niveaux (organigramme, categories
imbriquees, graphe de dependances).
\end{infobox}

% =============================================================================
\section{Fonctions fenetrees (window functions)}
% =============================================================================
Une fonction fenetree calcule une valeur \emph{a travers un ensemble de lignes
liees a la ligne courante} (la \emph{fenetre}), \textbf{sans} regrouper les
lignes comme le fait \texttt{GROUP BY} : chaque ligne du resultat reste
individuelle. C'est l'outil pour repondre a des questions comme « le rang de
chaque commande », « le cumul depuis le debut », ou « comparer une ligne a la
precedente ».

\begin{lstlisting}
-- Classement des commandes par montant, au sein de chaque client
SELECT client_id, montant,
       RANK() OVER (PARTITION BY client_id ORDER BY montant DESC) AS rang
FROM commandes;

-- Cumul progressif (running total) par client, trie par date
SELECT client_id, montant,
       SUM(montant) OVER (PARTITION BY client_id ORDER BY id) AS cumul
FROM commandes;

-- Comparer chaque ligne a la precedente / suivante
SELECT client_id, montant,
       LAG(montant)  OVER (PARTITION BY client_id ORDER BY id) AS montant_precedent,
       LEAD(montant) OVER (PARTITION BY client_id ORDER BY id) AS montant_suivant
FROM commandes;

-- ROW_NUMBER : numeroter les lignes -- utile pour dedupliquer
SELECT *
FROM (
    SELECT *, ROW_NUMBER() OVER (PARTITION BY email ORDER BY id) AS rn
    FROM clients
) t
WHERE rn = 1;    -- garde seulement la premiere ligne de chaque email en double
\end{lstlisting}

\begin{infobox}{RANK vs DENSE\_RANK vs ROW\_NUMBER}
Pour des valeurs ex aequo : \texttt{ROW\_NUMBER} attribue toujours un numero
unique (1, 2, 3...) sans egard aux egalites. \texttt{RANK} laisse un « trou »
apres une egalite (1, 2, 2, 4...). \texttt{DENSE\_RANK} n'en laisse pas (1, 2,
2, 3...).
\end{infobox}

% =============================================================================
\section{Types de jointure des ensembles}
% =============================================================================
\texttt{UNION}, \texttt{INTERSECT} et \texttt{EXCEPT} combinent les
\emph{resultats} de deux requetes ayant le meme nombre de colonnes et des
types compatibles -- a ne pas confondre avec les jointures, qui combinent des
\emph{colonnes}.

\begin{lstlisting}
-- UNION : combine et supprime les doublons ; UNION ALL les garde (plus rapide)
SELECT nom FROM clients_actifs
UNION
SELECT nom FROM clients_archives;

-- INTERSECT : lignes presentes dans les deux requetes
SELECT email FROM clients
INTERSECT
SELECT email FROM newsletter_abonnes;

-- EXCEPT : lignes de la premiere requete absentes de la seconde
SELECT email FROM clients
EXCEPT
SELECT email FROM newsletter_abonnes;
\end{lstlisting}

% =============================================================================
\section{Index et performance}
% =============================================================================
Un index est une structure auxiliaire (le plus souvent un arbre B) qui accelere
la recherche sur une colonne, au prix d'un espace disque supplementaire et d'un
cout a chaque ecriture (l'index doit etre mis a jour). Sans index, chercher une
valeur oblige a parcourir la table entiere (\emph{sequential scan}).

\begin{lstlisting}
-- Index simple sur une colonne souvent filtree
CREATE INDEX idx_commandes_client_id ON commandes(client_id);

-- Index composite : utile si les requetes filtrent sur les deux colonnes
-- ensemble (l'ordre des colonnes compte : la premiere doit etre la plus
-- selective ou la plus souvent utilisee seule)
CREATE INDEX idx_commandes_client_statut ON commandes(client_id, statut);

-- Verifier ce que le moteur fait reellement : lit-il l'index ou la table entiere ?
EXPLAIN ANALYZE
SELECT * FROM commandes WHERE client_id = 42;
\end{lstlisting}

\begin{infobox}{Quand indexer}
Indexer les colonnes utilisees dans \texttt{WHERE}, \texttt{JOIN ... ON} et
\texttt{ORDER BY} sur de grandes tables. Les cles primaires et (souvent) les
cles etrangeres sont deja indexees automatiquement par le SGBD (a verifier :
sur PostgreSQL, une cle etrangere n'est \textbf{pas} indexee automatiquement,
contrairement a la cle primaire).
\end{infobox}

\begin{piege}
Un index existe mais n'est pas utilise si la requete applique une fonction sur
la colonne indexee (\texttt{WHERE LOWER(email) = ...} n'utilise pas un index
simple sur \texttt{email}) ou si le filtre porte sur trop peu de lignes
distinctes (un index sur une colonne booleenne est rarement utile). Trop
d'index ralentissent aussi les ecritures (\texttt{INSERT}/\texttt{UPDATE}) : ce
n'est pas gratuit, il faut indexer avec intention.
\end{piege}

% =============================================================================
\section{Transactions}
% =============================================================================
Une transaction regroupe plusieurs operations en une seule unite
\emph{atomique} : soit toutes reussissent (\texttt{COMMIT}), soit aucune n'est
appliquee (\texttt{ROLLBACK}). Indispensable des qu'une operation logique
touche plusieurs tables -- par exemple, debiter un compte et en crediter un
autre ne doit jamais s'arreter a mi-chemin.

\begin{lstlisting}
BEGIN;

UPDATE comptes SET solde = solde - 100 WHERE id = 1;
UPDATE comptes SET solde = solde + 100 WHERE id = 2;

-- Si tout s'est bien passe :
COMMIT;

-- Si une verification echoue (ex. solde negatif) :
-- ROLLBACK;
\end{lstlisting}

\begin{infobox}{Proprietes ACID}
\textbf{A}tomicite (tout ou rien) -- \textbf{C}oherence (les contraintes
restent respectees) -- \textbf{I}solation (les transactions concurrentes ne se
perturbent pas) -- \textbf{D}urabilite (une fois validee, une transaction
survit meme a une panne).
\end{infobox}

% =============================================================================
\section{Erreurs et pieges frequents}
% =============================================================================
\renewcommand{\arraystretch}{1.35}
\begin{longtable}{>{\raggedright\arraybackslash}p{4.6cm}
                  >{\raggedright\arraybackslash}p{8.2cm}}
\rowcolor{primary}
\textbf{\color{white}Piege} & \textbf{\color{white}Explication} \\
\hline
\texttt{= NULL} au lieu de \texttt{IS NULL} & Ne renvoie jamais \texttt{TRUE} ;
utiliser \texttt{IS NULL} / \texttt{IS NOT NULL}. \\
\texttt{WHERE} sur une table jointe en \texttt{LEFT JOIN} & Transforme
silencieusement le \texttt{LEFT JOIN} en \texttt{INNER JOIN} ; filtrer dans le
\texttt{ON}. \\
Colonne non agregee absente du \texttt{GROUP BY} & Erreur (Postgres) ou valeur
arbitraire (MySQL en mode permissif). \\
\texttt{HAVING} vs \texttt{WHERE} & \texttt{WHERE} filtre avant regroupement,
\texttt{HAVING} apres. \\
\texttt{DELETE}/\texttt{UPDATE} sans \texttt{WHERE} & Affecte toutes les
lignes de la table. \\
\texttt{SELECT *} en production & Casse si les colonnes changent ; ramene des
donnees inutiles ; nommer les colonnes explicitement. \\
Comparer des types incompatibles & \texttt{'10' = 10} peut fonctionner par
conversion implicite mais degrade les performances (l'index n'est plus
utilisable) et masque des bugs. \\
Oublier les fuseaux horaires & Preferer \texttt{TIMESTAMPTZ} a
\texttt{TIMESTAMP} des que l'application a des utilisateurs dans plusieurs
fuseaux. \\
\end{longtable}

% =============================================================================
\section{Fiche recapitulative}
% =============================================================================
\renewcommand{\arraystretch}{1.35}
\begin{longtable}{>{\raggedright\arraybackslash}p{4.5cm}
                  >{\raggedright\arraybackslash}p{8.3cm}}
\rowcolor{primary}
\textbf{\color{white}Concept} & \textbf{\color{white}A retenir} \\
\hline
Ordre logique & \texttt{FROM > JOIN > WHERE > GROUP BY > HAVING > SELECT >
ORDER BY > LIMIT}. \\
\texttt{INNER} vs \texttt{LEFT JOIN} & \texttt{INNER} = correspondances
uniquement ; \texttt{LEFT} = tout de la table de gauche + \texttt{NULL}. \\
\texttt{WHERE} vs \texttt{HAVING} & Avant regroupement vs apres regroupement. \\
\texttt{IN} vs \texttt{EXISTS} & \texttt{EXISTS} s'arrete des la premiere
correspondance, souvent plus rapide sur de gros volumes. \\
CTE (\texttt{WITH}) & Nomme une sous-requete ; \texttt{WITH RECURSIVE} pour les
hierarchies. \\
Fonctions fenetrees & \texttt{OVER (PARTITION BY ... ORDER BY ...)} : calcule
sans regrouper les lignes. \\
Index & Accelere la lecture, ralentit l'ecriture ; verifier avec
\texttt{EXPLAIN ANALYZE}. \\
Transaction & \texttt{BEGIN / COMMIT / ROLLBACK} ; proprietes ACID. \\
\texttt{NULL} & Valeur inconnue ; toujours \texttt{IS NULL}, jamais \texttt{=
NULL}. \\
\end{longtable}

\vspace{0.6cm}
\begin{infobox}{Methode pour ecrire une requete complexe}
\begin{enumerate}[leftmargin=*,itemsep=2pt,topsep=2pt]
  \item Identifier les tables necessaires et comment elles se relient (cles).
  \item Ecrire d'abord le \texttt{FROM}/\texttt{JOIN}, verifier avec
  \texttt{SELECT *} que les lignes obtenues sont les bonnes.
  \item Ajouter les filtres (\texttt{WHERE}), puis le regroupement
  (\texttt{GROUP BY}/\texttt{HAVING}) si necessaire.
  \item Ne selectionner les colonnes finales et trier (\texttt{ORDER BY}) qu'en
  dernier.
  \item Sur une grosse table, verifier le plan d'execution
  (\texttt{EXPLAIN ANALYZE}) avant de considerer la requete terminee.
\end{enumerate}
\end{infobox}

\end{document}
